% Auto-generated by Code2Latex
% Requires: longtable, booktabs packages

\begin{longtable}{p{\textwidth}}
\caption{Tools for single (Level 1)} \\
\endfirsthead

\multicolumn{1}{c}{\tablename\ \thetable{} -- continued from previous page} \\
\endhead

\multicolumn{1}{r}{Continued on next page} \\
\endfoot

\endlastfoot

\textbf{\texttt{batch\_retrieve\_polymorphs}} \\
\begin{description}
    \item[\textbf{Description:}] Retrieve polymorphs for multiple chemical compositions efficiently in batch mode and save it to given directory as json.

This tool performs polymorph retrieval for multiple chemical compositions simultaneously.
The tool applies energy and count filters to focus on thermodynamically relevant phases.
This tool retrieves comprehensive polymorph data from the Materials Project database for each chemical composition and save in the user inputted save\_directory.
Polymorphs are different crystal structures with the same chemical formula but different atomic arrangements, leading to distinct physical and chemical properties.
This tool could be relevant for retrieving structures of the same composition for multiple composition.
Apart from structure for each polymorph Materials Project ID (MP ID), CIF structure, energy above hull, formation energy per atom, band gap, density, volume, number of sites, space group, and stability information is also retrieved.
The result is then saved a JSON file at the specified directory in the format <composition>\_polymorphs.json.
    \item[\textbf{Arguments:}]
    \begin{description}
        \item \texttt{compositions} (list[str]): List of chemical compositions to retrieve.
List of chemical formulas for which polymorphs should be retrieved. Each composition should follow standard chemical notation. The tool will process each composition independently and provide detailed success/failure reporting. Large lists are supported but may take significant time to process.
        \item \texttt{max\_energy\_above\_hull} (float): Maximum energy above hull threshold in eV/atom. Defaults to 0.5.
Energy threshold above the convex hull for including polymorphs. Only phases with energy above hull less than or equal to this value will be included. This filters out highly unstable phases while retaining potentially accessible metastable phases. Lower values give more stable phases but may miss interesting metastable structures.
        \item \texttt{max\_per\_composition} (int): Maximum number of polymorphs per composition. Defaults to 10.
Maximum number of polymorphs to retrieve for each composition, taken from the most stable phases first. This prevents data explosion for compositions with many known phases while ensuring the most important structures are captured. Higher values provide more comprehensive coverage but increase dataset size and processing time.
        \item \texttt{save\_directory} (str): Directory path for saving individual composition files. Defaults to "polymorph\_data".
Base directory where individual JSON files for each composition will be saved. The directory will be created if it doesn't exist. Each composition will have its own JSON file named with the composition formula. This organization facilitates easy data management and selective loading of specific compositions.
    \end{description}
    \item[\textbf{Returns:}] JSON string with batch retrieval results and statistics.

A comprehensive JSON report containing lists of successfully processed and failed compositions, total number of polymorphs retrieved, file paths for each composition, and summary statistics.
         This enables quality control and tracking of the batch processing workflow.
\end{description} \\
\midrule
\textbf{\texttt{cat\_files}} \\
\begin{description}
    \item[\textbf{Description:}] Concatenate and display contents of one or more files
    \item[\textbf{Arguments:}]
    \begin{description}
        \item \texttt{paths} (list[str]): List of file paths to concatenate
        \item \texttt{separator} (str): Separator between file contents
    \end{description}
\end{description} \\
\midrule
\textbf{\texttt{consolidate\_polymorph\_datasets}} \\
\begin{description}
    \item[\textbf{Description:}] Consolidate multiple polymorph JSON files into a single comprehensive dataset.

This tool combines multiple polymorph datasets from different compositions into a unified dataset suitable for dataset preparation and machine learning applications.
It handles data integration and provides comprehensive statistics about the consolidated dataset.
Often you have multiple polymorph json file and you want to combine them. You can use this tool to combine them
    \item[\textbf{Arguments:}]
    \begin{description}
        \item \texttt{composition\_files} (dict[str, str]): Dictionary mapping compositions to their JSON file paths.
A dictionary where keys are composition names/formulas and values are file paths to their corresponding JSON files containing polymorph data. The tool will attempt to read each file and integrate the data while maintaining composition information.
        \item \texttt{output\_path} (str): Path for the consolidated dataset file. Defaults to "consolidated\_polymorphs.json".
File path where the consolidated dataset will be saved. The file will contain all polymorphs from all compositions in a single JSON structure with added source composition information. The directory will be created if it doesn't exist. Using .json extension is recommended for clarity.
    \end{description}
    \item[\textbf{Returns:}] JSON string with consolidation results and comprehensive statistics.

A JSON-formatted string containing consolidation status, output file path, and detailed statistics including total number of polymorphs, number of compositions successfully included, average polymorphs per composition, and any processing errors.
         This enables quality control and assessment of the consolidation process.
\end{description} \\
\midrule
\textbf{\texttt{copy\_file}} \\
\begin{description}
    \item[\textbf{Description:}] Copy a file from a source to destination
    \item[\textbf{Arguments:}]
    \begin{description}
        \item \texttt{source} (str): Source file path
        \item \texttt{destination} (str): Destination file path
    \end{description}
\end{description} \\
\midrule
\textbf{\texttt{evaluate\_xgboost\_model}} \\
\begin{description}
    \item[\textbf{Description:}] Evaluate trained XGBoost model with comprehensive performance metrics and detailed analysis.

This tool provides thorough evaluation of trained XGBoost models with comprehensive metrics and detailed analysis capabilities.
It generates standard regression metrics (MAE, RMSE, R2, MAPE), error analysis (mean error, error standard deviation, max positive/negative errors), prediction ranges (min, max, standard deviation), and top feature importance rankings essential for model validation and deployment decisions.
The tool supports both basic and detailed analysis modes, enabling quick assessments or in-depth model understanding for research and production applications.
    \item[\textbf{Arguments:}]
    \begin{description}
        \item \texttt{model\_path} (str): Path to the saved XGBoost model file.
Complete file path to the serialized XGBoost model. The model should be saved using joblib or pickle and contain a trained XGBoost regressor ready for evaluation. The file must be readable and contain a valid model object.
        \item \texttt{test\_data\_path} (str): Path to test data CSV file with same structure as training data.
Complete file path to CSV file containing test data with identical column structure to the training data used for model creation. Must include both feature columns and the target column for evaluation. The data should be preprocessed consistently with the training data.
        \item \texttt{target\_column} (str): Name of the target column for evaluation. Defaults to "formation\_energy\_per\_atom".
The column name in the test CSV that contains the true values for comparison with model predictions. This should match the target column used during training. Common targets include formation energy, band gap, and other materials properties.
        \item \texttt{detailed\_analysis} (bool): Whether to include detailed analysis and feature importance. Defaults to True.
Boolean flag controlling the depth of analysis performed. When True, includes prediction ranges (min, max, std), error analysis (mean error, error std, max errors), and top 10 feature importance rankings. When False, provides only basic metrics (MAE, RMSE, R2, MAPE) for quick assessment. Detailed analysis is recommended for thorough model evaluation and interpretation.
    \end{description}
    \item[\textbf{Returns:}] JSON string with comprehensive evaluation metrics and analysis results.

A detailed JSON-formatted string containing evaluation success status, comprehensive performance metrics (MAE, RMSE, R2, MAPE), prediction statistics (min/max/std), error analysis (mean error, error std, max positive/negative errors), and top 10 feature importance rankings when detailed analysis is enabled.
         This provides complete model assessment for validation and comparison purposes.
\end{description} \\
\midrule
\textbf{\texttt{file\_info}} \\
\begin{description}
    \item[\textbf{Description:}] Get information about a file or directory
    \item[\textbf{Arguments:}]
    \begin{description}
        \item \texttt{path} (str): Path to the file or directory
    \end{description}
\end{description} \\
\midrule
\textbf{\texttt{filter\_json\_with\_strategy}} \\
\begin{description}
    \item[\textbf{Description:}] Filter JSON data using custom Python code and save results to a new file.

This tool provides flexible JSON data filtering capabilities using custom Python logic, essential for data preprocessing, quality control, and custom analysis workflows.
It enables sophisticated filtering operations that go beyond simple threshold-based selection, allowing for complex multi-criteria filtering, data validation, and custom transformations.
This is crucial for preparing datasets for analysis and machine learning applications.
    \item[\textbf{Arguments:}]
    \begin{description}
        \item \texttt{input\_json\_path} (str): Path to the input JSON file to be filtered.
Complete file path to the JSON file containing the data to be filtered. The file should contain valid JSON data, typically a list of dictionaries representing materials or structures with various properties. The file must be readable and contain well-formed JSON.
        \item \texttt{output\_json\_path} (str): Path where filtered JSON data will be saved.
Complete file path where the filtered JSON data will be written. The directory will be created if it doesn't exist. The output file will contain the filtered subset of the input data in the same JSON format. Using .json extension is recommended for clarity.
        \item \texttt{custom\_code} (str | None): Python code string defining the filtering logic.
A string containing Python code that defines the filtering logic. The code should expect the input data in a variable named 'data' and store the filtered results in a variable named 'filtered\_data'. The code can use any Python constructs including list comprehensions, complex conditions, and data transformations.
    \end{description}
    \item[\textbf{Returns:}] JSON string with filtering results and comprehensive statistics.

A JSON-formatted string containing filtering status, original and filtered data counts, output file path, and percentage reduction achieved.
         This provides comprehensive information about the filtering operation's effectiveness and enables quality control of the data processing pipeline.
\end{description} \\
\midrule
\textbf{\texttt{get\_bulk\_polymorphs\_data}} \\
\begin{description}
    \item[\textbf{Description:}] Query Materials Project database to find all polymorphs for a given chemical composition.

This tool retrieves comprehensive polymorph data from the Materials Project database for a specific chemical composition.
Polymorphs are different crystal structures with the same chemical formula but different atomic arrangements, leading to distinct physical and chemical properties.
This tool could be relevant for retrieving structures of the same composition.
Apart from structure for each polymorph Materials Project ID (MP ID), CIF structure, energy above hull, formation energy per atom, band gap, density, volume, number of sites, space group, and stability information is also retrieved.
    \item[\textbf{Arguments:}]
    \begin{description}
        \item \texttt{composition} (str): Chemical composition formula.
Chemical formula specifying the composition for which polymorphs should be retrieved. Should follow standard chemical notation with element symbols. The tool will find all known crystal structures with this exact composition in the Materials Project database.
    \end{description}
    \item[\textbf{Returns:}] JSON string containing comprehensive polymorph data sorted by stability.

A JSON-formatted string containing a list of dictionaries, each representing a polymorph with properties including Materials Project ID, CIF structure, energy above hull, formation energy per atom, band gap, density, volume, number of sites, space group, and stability information.
         Results are sorted by energy above hull for easy identification of the most stable phases.
\end{description} \\
\midrule
\textbf{\texttt{get\_bulk\_polymorphs\_data\_to\_file}} \\
\begin{description}
    \item[\textbf{Description:}] Query Materials Project for polymorphs and save comprehensive data to a JSON file to give path.

This tool performs the same comprehensive polymorph retrieval  but saves the results directly to a JSON file for persistent storage and later analysis.
This tool retrieves comprehensive polymorph data from the Materials Project database for a specific chemical composition.
Polymorphs are different crystal structures with the same chemical formula but different atomic arrangements, leading to distinct physical and chemical properties.
This tool could be relevant for retrieving structures of the same composition.
Apart from structure for each polymorph Materials Project ID (MP ID), CIF structure, energy above hull, formation energy per atom, band gap, density, volume, number of sites, space group, and stability information is also retrieved.
The result is then saved a JSON file at the specified save\_path.
The file-based approach allows for efficient handling of large datasets and facilitates saving context of llm.
    \item[\textbf{Arguments:}]
    \begin{description}
        \item \texttt{composition} (str): Chemical composition formula.
Chemical formula specifying the composition for which polymorphs should be retrieved and saved. Should follow standard chemical notation with element symbols and subscripts. The tool will find all known crystal structures with this exact composition in the Materials Project database.
        \item \texttt{save\_path} (str | None): File path where JSON data will be saved.
Complete file path including filename and extension where the polymorph data will be saved. The path should be writable and the directory will be created if it doesn't exist. Using .json extension is recommended for clarity. If None, the tool will raise an error as the file path is required.
    \end{description}
    \item[\textbf{Returns:}] File path where the polymorph data was saved.

Returns the exact file path where the JSON data was successfully written.
         This path can be used by subsequent tools for data loading and processing.
         The file contains comprehensive polymorph data in JSON format, sorted by thermodynamic stability.
\end{description} \\
\midrule
\textbf{\texttt{list\_files}} \\
\begin{description}
    \item[\textbf{Description:}] List files in a directory
    \item[\textbf{Arguments:}]
    \begin{description}
        \item \texttt{path} (str): Path to the directory
        \item \texttt{recursive} (bool): Whether to list files recursively
    \end{description}
\end{description} \\
\midrule
\textbf{\texttt{perform\_cross\_validation}} \\
\begin{description}
    \item[\textbf{Description:}] Perform k-fold cross-validation to assess model stability and generalization performance.

This tool implements comprehensive k-fold cross-validation for XGBoost models to assess model generalization capability, and robustness across different data splits.
Cross-validation provides more reliable performance estimates than single train/test splits by evaluating model performance across multiple data partitions.
This is essential for hyperparameter tuning, model comparison, and ensuring reliable performance estimates for materials property prediction models.
    \item[\textbf{Arguments:}]
    \begin{description}
        \item \texttt{train\_data\_path} (str): Path to training data CSV file.
Complete file path to CSV file containing training data with features and target column. The file should have a header row and be properly formatted for machine learning. This data will be split into k folds for cross-validation, so it should represent the complete training dataset.
        \item \texttt{target\_column} (str): Name of the target column for prediction. Defaults to "formation\_energy\_per\_atom".
The column name in the CSV that contains the target values for prediction. This column will be separated from features during cross-validation. Should match the target used in subsequent training workflows. Common targets include formation energy, band gap, and other materials properties.
        \item \texttt{cv\_folds} (int): Number of cross-validation folds. Defaults to 5.
The number of folds to use for k-fold cross-validation. Higher values provide more robust estimates but increase computational cost. Common choices are 5 or 10 folds. The value should be chosen based on dataset size - smaller datasets benefit from more folds while larger datasets can use fewer folds.
        \item \texttt{hyperparameters} (dict | None): Optional XGBoost hyperparameters for cross-validation.
Dictionary containing XGBoost hyperparameters to use across all cross-validation folds. If None, optimized default parameters will be used. Consistent hyperparameters across folds ensure fair comparison and reliable performance estimates. Useful for testing specific hyperparameter configurations.
    \end{description}
    \item[\textbf{Returns:}] JSON string with comprehensive cross-validation results and statistical analysis.

A detailed JSON-formatted string containing cross-validation success status, individual fold scores, statistical summaries (mean, standard deviation), hyperparameters used, and performance stability assessment.
         This enables comprehensive evaluation of model robustness and generalization capability.
\end{description} \\
\midrule
\textbf{\texttt{prepare\_tabular\_dataset}} \\
\begin{description}
    \item[\textbf{Description:}] Prepare tabular dataset for traditional ML models with feature engineering.

This tool creates ML-ready tabular datasets from materials data with simple feature engineering capabilities suitable for traditional machine learning models like XGBoost.
It implements multiple feature engineering strategies such as basic property extraction to structural descriptors, handles data preprocessing, normalization, and train/test splitting.
This is essential for building property prediction models.
    \item[\textbf{Arguments:}]
    \begin{description}
        \item \texttt{polymorphs\_json\_path} (str): Path to consolidated polymorphs JSON file.
Complete file path to a JSON file containing consolidated polymorph data with materials properties and crystal structures.
        \item \texttt{output\_path} (str): Base path for saving dataset files.
Base directory and filename prefix where the prepared dataset files will be saved. Multiple files will be created including training data, test data, and metadata. The tool will create the directory structure if it doesn't exist.
        \item \texttt{target\_property} (str): Property to predict. Defaults to "formation\_energy\_per\_atom".
The materials property that will serve as the prediction target for machine learning models. This should be a key present in the polymorphs data with numerical values. Common targets include formation energy, band gap, bulk modulus, and other calculated properties.
        \item \texttt{feature\_engineering} (str): Feature engineering strategy. Defaults to "basic".
The level of feature engineering to apply to the materials data. "basic" extracts fundamental properties, "advanced" includes structural and electronic descriptors, and "custom" applies specialized feature extraction.
        \item \texttt{test\_split} (float): Fraction of data for test set. Defaults to 0.2.
The proportion of the dataset to reserve for testing, expressed as a decimal fraction. The remaining data will be used for training. Common values range from 0.1 to 0.3 depending on dataset size and validation strategy. Larger test sets provide more reliable evaluation but reduce training data.
        \item \texttt{normalize} (bool): Whether to normalize features. Defaults to True.
Boolean flag controlling whether features should be normalized using standard scaling (zero mean, unit variance). Normalization is generally recommended for most ML algorithms as it ensures features have similar scales and prevents any single feature from dominating the model.
    \end{description}
    \item[\textbf{Returns:}] JSON string with comprehensive dataset preparation results and file paths.

A detailed JSON-formatted string containing preparation success status, file paths for training and test data, normalization parameters, dataset statistics, feature information, and metadata.
         This provides complete information about the prepared dataset for subsequent ML workflows.
\end{description} \\
\midrule
\textbf{\texttt{read\_file}} \\
\begin{description}
    \item[\textbf{Description:}] Read the contents of a file into a string
    \item[\textbf{Arguments:}]
    \begin{description}
        \item \texttt{path} (str): Path to the file to read
    \end{description}
\end{description} \\
\midrule
\textbf{\texttt{select\_polymorphs\_with\_strategy}} \\
\begin{description}
    \item[\textbf{Description:}] Select polymorphs using strategic criteria for systematic materials analysis.

This tool implements intelligent selection strategies for polymorph datasets, enabling systematic reduction of large materials databases while preserving important structural and energetic diversity.
The input can be json polymorph data as string or path to json file of polymorph data.
It supports multiple selection algorithms designed for different research objectives, from stability-focused studies to comprehensive structural surveys.
This is crucial for managing computational resources and focusing analysis on the most relevant materials.
    \item[\textbf{Arguments:}]
    \begin{description}
        \item \texttt{polymorphs\_data} (str): JSON string or file path containing polymorph data.
Either a JSON-formatted string containing polymorph data or a file path to a JSON file, depending on the is\_path parameter. The data should contain polymorphs with properties like energy\_above\_hull, space\_group, and other structural/energetic information. This is typically output from polymorph retrieval tools.
        \item \texttt{selection\_strategy} (str): Strategy for polymorph selection. Defaults to "diverse\_energy".
The algorithm used for selecting polymorphs from the dataset "diverse\_energy" selects polymorphs distributed across the energy range for representative sampling. "most\_stable" prioritizes the most thermodynamically stable phases. "diverse\_structure" ensures different space groups are represented to capture structural diversity.
        \item \texttt{max\_polymorphs} (int): Maximum number of polymorphs to select. Defaults to 5.
The maximum number of polymorphs to include in the final selection. This parameter controls the size of the resulting dataset and should be chosen based on computational resources and analysis requirements. Larger values provide more comprehensive coverage but increase processing time and computational cost.
        \item \texttt{energy\_threshold} (float): Maximum energy above hull in eV/atom. Defaults to 0.5.
Energy threshold above the convex hull for including polymorphs in the selection process. Only phases with energy above hull less than or equal to this value will be considered. This pre-filtering step ensures that only thermodynamically accessible phases are included in the analysis.
        \item \texttt{is\_path} (bool): Whether polymorphs\_data is a file path. Defaults to False.
Boolean flag indicating whether the polymorphs\_data parameter should be treated as a file path (True) or as a JSON string (False). When True, the tool will read the JSON data from the specified file. When False, it will parse the data directly from the string.
    \end{description}
    \item[\textbf{Returns:}] JSON string containing selected polymorphs based on the specified strategy.

A JSON-formatted string containing the selected subset of polymorphs, maintaining the same data structure as the input but with reduced number of entries.
         The selection preserves important characteristics according to the chosen strategy while reducing dataset size for efficient processing.
\end{description} \\
\midrule
\textbf{\texttt{select\_polymorphs\_with\_strategy\_to\_file}} \\
\begin{description}
    \item[\textbf{Description:}] Select polymorphs using strategic criteria and save results to file for persistent storage.

This tool combines the strategic polymorph selection capabilities with direct file output for persistent storage and workflow automation.
It implements the same selection algorithms as select\_polymorphs\_with\_strategy but automatically saves results to a specified file path.
The input can be json polymorph data as string or path to json file of polymorph data.
It supports multiple selection algorithms designed for different research objectives, from stability-focused structures to comprehensive structures.
This is essential for automated workflows, batch processing, and creating organized datasets where selected polymorphs need to be stored for later use or sharing.
    \item[\textbf{Arguments:}]
    \begin{description}
        \item \texttt{polymorphs\_data} (str): JSON string or file path containing polymorph data.
Either a JSON-formatted string containing polymorph data or a file path to a JSON file, depending on the is\_path parameter. The data should contain polymorphs with properties like energy\_above\_hull, space\_group, and other structural/energetic information for strategic selection.
        \item \texttt{save\_path} (str): File path where selected polymorphs will be saved.
Complete file path where the selected polymorph subset will be saved in JSON format. The directory will be created if it doesn't exist. This file can be used by subsequent tools or shared with collaborators. Using .json extension is recommended for clarity.
        \item \texttt{selection\_strategy} (str): Strategy for polymorph selection. Defaults to "diverse\_energy".
The algorithm used for selecting polymorphs from the dataset. Options include "diverse\_energy" for energy range sampling, "most\_stable" for thermodynamic stability, and "diverse\_structure" for structural diversity. Each strategy optimizes for different research objectives and analysis requirements.
        \item \texttt{max\_polymorphs} (int): Maximum number of polymorphs to select. Defaults to 5.
The maximum number of polymorphs to include in the final selection and save to file. This parameter controls dataset size and should be chosen based on computational resources and analysis requirements. Larger values provide more comprehensive coverage but increase processing time.
        \item \texttt{energy\_threshold} (float): Maximum energy above hull in eV/atom. Defaults to 0.5.
Energy threshold above the convex hull for including polymorphs in the selection process. Only phases with energy above hull less than or equal to this value will be considered. This pre-filtering ensures thermodynamic accessibility of selected phases.
        \item \texttt{is\_path} (bool): Whether polymorphs\_data is a file path. Defaults to False.
Boolean flag indicating whether the polymorphs\_data parameter should be treated as a file path (True) or as a JSON string (False). When True, the tool will read the JSON data from the specified file. This enables flexible input handling for different workflow patterns.
    \end{description}
    \item[\textbf{Returns:}] File path where the selected polymorphs were saved.

The complete file path where the selected polymorphs have been successfully saved.
         This path can be used by subsequent tools for loading the selected dataset or for verification that the file was created correctly.
         The file contains the subset of polymorphs selected according to the specified strategy.
\end{description} \\
\midrule
\textbf{\texttt{sort\_and\_get\_first\_from\_json}} \\
\begin{description}
    \item[\textbf{Description:}] Sort JSON data by specified key and return the first element's specified value.

This utility tool provides flexible sorting and extraction capabilities for JSON data, particularly useful for materials data analysis where you need to identify optimal structures based on specific criteria. It enables quick identification of the best material according to any numerical property, such as finding the most stable phase, highest band gap material, or densest structure.
    \item[\textbf{Arguments:}]
    \begin{description}
        \item \texttt{polymorph\_data\_json} (str): JSON string containing the data to be sorted.
A JSON-formatted string containing a list of dictionaries, each representing a material or structure with various properties. The data should be structured consistently with numerical values for the sorting key. This is typically output from polymorph retrieval tools.
        \item \texttt{sort\_key} (str): Property name to sort the data by.
The dictionary key name that will be used for sorting the data. This should correspond to a numerical property in the JSON data. The sorting is performed in ascending order, so the first element will have the smallest value for this property. Common keys include energy\_above\_hull, band\_gap, density, formation\_energy\_per\_atom.
        \item \texttt{return\_key} (str): Property name to return from the first element after sorting.
The dictionary key name for the value that should be returned from the first (optimal) element after sorting. This allows extraction of any property from the optimal structure, such as material\_id for identification, cif for structure, or any other calculated property.
    \end{description}
    \item[\textbf{Returns:}] Value of the specified return\_key from the first element after sorting.

The value corresponding to the return\_key from the material that has the smallest value for the sort\_key. This could be a string (like material\_id or CIF), a number (like energy or band gap), or any other data type stored in the JSON. The returned value represents the optimal material according to the specified sorting criterion.
\end{description} \\
\midrule
\textbf{\texttt{train\_xgboost\_model}} \\
\begin{description}
    \item[\textbf{Description:}] Train XGBoost regression model for property prediction with evaluation.

This tool implements comprehensive XGBoost model training for materials property prediction,including hyperparameter management, model evaluation.
The tool provides complete training pipeline with automatic evaluation metrics and save the model to the give path.
The model takes as input csv file for train and test dataset
    \item[\textbf{Arguments:}]
    \begin{description}
        \item \texttt{train\_data\_path} (str): Path to training data CSV file.
Complete file path to the CSV file containing training data with features and target column. The file should have a header row with column names and be properly formatted with numerical features.
        \item \texttt{test\_data\_path} (str): Path to test data CSV file.
Complete file path to the CSV file containing test data with the same structure as training data. Used for independent model evaluation and performance assessment. Should have identical column structure to training data.
        \item \texttt{model\_save\_path} (str): Path to save the trained model file in .pkl format
Complete file path where the trained XGBoost model will be saved using joblib serialization. The model can be loaded later for predictions or further analysis. Using .pkl extension is recommended for clarity.
        \item \texttt{target\_column} (str): Name of the target column for prediction. Defaults to "formation\_energy\_per\_atom".
The column name in the CSV files that contains the target values to predict. This column will be separated from features during training. Common targets include formation energy, band gap, bulk modulus, and other materials properties.
        \item \texttt{hyperparameters} (dict | None): Optional dictionary of XGBoost hyperparameters.
Dictionary containing XGBoost hyperparameters to override default values. Can include parameters like n\_estimators, max\_depth, learning\_rate, subsample, etc. If None, optimized default parameters will be used. Proper hyperparameter tuning can significantly improve model performance.
    \end{description}
    \item[\textbf{Returns:}] JSON string with comprehensive training results and model performance metrics.

A detailed JSON-formatted string containing training success status, model performance metrics (MAE, RMSE, R2), feature importance rankings, hyperparameters used, dataset information, and file paths for saved model and results.
         This enables comprehensive model evaluation and comparison.
\end{description} \\
\midrule
\textbf{\texttt{write\_file}} \\
\begin{description}
    \item[\textbf{Description:}] Write content to a file
    \item[\textbf{Arguments:}]
    \begin{description}
        \item \texttt{path} (str): Path to the file to write
        \item \texttt{content} (str): Content to write to the file
    \end{description}
\end{description} \\
\end{longtable}
